% ---------------------------------------------------------------------------
%  PART III -- The question papers, reproduced in full
% ---------------------------------------------------------------------------
\chapter{The Question Papers}
\chaptermark{The Question Papers}

The five papers for which full scans survive are reproduced here
question-for-question, in chronological order, exactly as printed --- header,
instructions, question order, numbering, and the three code columns each paper
carried beside its questions.

\vspace{2pt}
Every question also carries a pointer in the last column to the answer
elsewhere in this book, so a whole past paper can be worked end to end and
then marked:

\begin{keybox}
\textbf{Q\emph{n}} \ $\rightarrow$ \ Part~I, descriptive question \emph{n}
\hfill
\textbf{A\emph{n}} \ $\rightarrow$ \ Part~II, multiple-choice item \emph{n}
\newline
\textbf{B\emph{n}} \ $\rightarrow$ \ Part~II, fill-in-the-blank \emph{n}
\hfill
\textbf{C\emph{n}} \ $\rightarrow$ \ Part~II, one-mark short answer \emph{n}
\end{keybox}

The four papers not reproduced here --- SEE October 2023, CIE~1 of 22 June
2024, CIE~2 of 24 July 2024 and SEE September/October 2024 --- are listed
question by question in \textbf{Appendix~A}, with the same pointers. Their
questions are all answered; only the printed paper itself was unavailable.

\section{Improvement CIE --- 28 August 2024}

\begin{papermeta}
\mk{Course code}   & \mv{HS248AT}                    &
\mk{Maximum marks} & \mv{25}                          \\
\mk{Semester}      & \mv{IV}                          &
\mk{Duration}      & \mv{60 minutes}                  \\
\mk{Academic year} & \mv{2023--2024, Even Semester}   &
\mk{Assessment}    & \mv{Improvement CIE --- Quiz and Test} \\
\mk{Date}          & \mv{28 August 2024}              &
\mk{Structure}     & \mv{Part A 5 M \textbf{+} Part B 25 M} \\
\end{papermeta}

\subsection*{Part A --- Objective Questions \hfill 5 Marks}

\begin{qtable}
\qn{1} & What is the fundamental principle of harmony in nature and existence?
  \optsinline{Competition}{Coexistence}{Dominance}{Isolation}
  & 1 & CO3 & L1 & \ap{A1} \\ \hline
\qn{2} & What is the role of human beings in maintaining harmony in nature?
  \begin{optlist}
  \item Exploiting natural resources
  \item Preserving and nurturing nature
  \item Dominating other species
  \item Isolating from the environment
  \end{optlist}
  & 1 & CO2 & L2 & \ap{A2} \\ \hline
\qn{3} & Which of the following best describes the Physical Order?
  \begin{optlist}
  \item Includes only human-made objects
  \item Comprises natural elements like soil, water, air, and minerals
  \item Refers exclusively to living beings
  \item Focuses on mental and spiritual aspects
  \end{optlist}
  & 1 & CO2 & L1 & \ap{A3} \\ \hline
\qn{4} & What distinguishes the Human Order from the other three orders?
  \begin{optlist}
  \item Physical strength
  \item Intelligence and moral values
  \item Ability to photosynthesize
  \item Instinctual behaviour only
  \end{optlist}
  & 1 & CO4 & L1 & \ap{A4} \\ \hline
\qn{5} & In the context of coexistence, what is the meaning of `prosperity'?
  \begin{optlist}
  \item Material wealth accumulation
  \item Holistic well-being and happiness
  \item Physical strength and power
  \item Technological superiority
  \end{optlist}
  & 1 & CO3 & L1 & \ap{A5} \\ \hline
\end{qtable}

\subsection*{Part B --- Descriptive Questions \hfill 25 Marks}

\begin{qtable}
\qn{1} & Existence is co-existence of mutually interacting units in
  all-pervasive space. Explain & 5 & CO3 & L2 & \ap{Q62} \\ \hline
\qn{2} & What is the natural characteristics (swabhava) of human order?
  Explain & 5 & CO2 & L1 & \ap{Q54} \\ \hline
\qn{3} & What are the four orders of nature? Briefly explain them
  & 5 & CO4 & L2 & \ap{Q51} \\ \hline
\qn{4} & How can we say that `nature is Self-Organized and in space
  Self-Organization Is Available.' & 5 & CO4 & L2 & \ap{Q61} \\ \hline
\qn{5} & Explain the difference and similarities between Pranic order and
  animal order. What is the relation between the two orders? & 5 & CO2 & L1 &
  \ap{Q57} \\ \hline
\end{qtable}

\subsection*{Marks distribution as printed}

\begin{mdtable}
Test --- Max Marks & xx & 12 & 07 & 11 & 14 & 16 & xx & xx & xx & xx \\ \hline
\end{mdtable}

\vspace{4pt}
Four cells were left as the template's own \texttt{xx} placeholder and are
reproduced that way. They are all cells that would have been zero: no question
on this paper is mapped to CO1, and none is set above L2.

\begin{note}
\textbf{\sffamily The paper does not add up to its own maximum.} The header
states a maximum of \textbf{25 marks}, but Part~A carries 5 and the five
Part~B questions carry 5 each --- \textbf{30 in total}. The footer table above
agrees with the questions, not with the header: its CO row totals 30
($12+7+11$) and its BT row totals 30 ($14+16$), and both reconcile exactly
with the per-question codes in the two grids. The discrepancy is in the
original and has been left as printed.
\end{note}

\section{CIE-II --- 7 May 2025}

\begin{papermeta}
\mk{Course code}   & \mv{HS248AT}                    &
\mk{Maximum marks} & \mv{25 $+$ 5}                    \\
\mk{Semester}      & \mv{IV}                          &
\mk{Duration}      & \mv{60 minutes}                  \\
\mk{Academic year} & \mv{2024--2025, Even Semester}   &
\mk{Assessment}    & \mv{CIE-II --- Quiz and Test}    \\
\mk{Date}          & \mv{7 May 2025}                  &
\mk{Structure}     & \mv{Part A 5 M \textbf{+} Part B 25 M} \\
\end{papermeta}

This is the only regular CIE in the collection --- the other four papers are
improvement or semester-end sittings. Its subject matter is correspondingly
narrower: every question on it comes from harmony in the family and the values
in relationship, which is the block of the syllabus a second CIE covers.

\subsection*{Part A --- Objective Questions \hfill 5 Marks}

\begin{qtable}
\qn{1} & Which of the following is considered the foundation of a strong human
  relationship?
  \optsinline{Intelligence}{Trust}{Wealth}{Appearance}
  & 1 & CO1 & L1 & \ap{A17} \\ \hline
\qn{2} & Disrespect Arising out of
  \begin{optlist}
  \item Lack of communication
  \item Lack of knowledge
  \item Differentiation leading to Discrimination
  \item Lack of Wealth
  \end{optlist}
  & 1 & CO2 & L1 & \ap{A18} \\ \hline
\qn{3} & Care is the feeling
  \begin{optlist}
  \item of responsibility and commitment for nurturing and protection of the
        Body of my relative.
  \item of responsibility and providing facility for nurturing and protection
        of the Body of my relative
  \item of responsibility of making the other feel comfortable with respect to
        Body of my relative
  \item of showing affection to the Body of my relative
  \end{optlist}
  & 1 & CO2 & L1 & \ap{A19} \\ \hline
\qn{4} & Fulfilment of relationship means
  \begin{optlist}
  \item Making effort for self development, i.e. development of one's own
        competence
  \item Making effort for mutual development, i.e. development of one's own
        competence and being of help to the other in developing their
        competence
  \item Making effort for mutual understanding, i.e. understanding of one's
        own competence and being of help to the other in understanding their
        competence
  \item Making effort for mutual in acquiring prosperity
  \end{optlist}
  & 1 & CO3 & L1 & \ap{A20} \\ \hline
\qn{5} & With the complete understanding of respect
  \begin{optlist}
  \item we can see for every individual on the earth that we all are the same
        in terms of intelligent, prosperity and status.
  \item we can see for every individual on the earth that we all are the same
        in terms of Knowledge, family, society.
  \item we can see for every individual on the earth that we all are the same
        in terms of competence, prosperity and status.
  \item we can see for every individual on the earth that we all are the same
        in terms of intention, program and potential.
  \end{optlist}
  & 1 & CO1 & L1 & \ap{A21} \\ \hline
\end{qtable}

\subsection*{Part A --- marks distribution as printed}

\begin{mdtable}
Test --- Max Marks & 2 & 2 & 1 & --- & 5 & --- & --- & --- & --- & --- \\
\hline
\end{mdtable}

\subsection*{Part B --- Descriptive Questions \hfill 25 Marks}

\begin{qtable}
\qn{1} & Justify the following (i)~Relation ship is between one Self (I1) and
  another Self (I2), (ii)~feelings in relationship --- in one Self (I1) for the
  other Self (I2) & 2$+$3 & CO1 & L1 & \ap{Q27} \\ \hline
\qn{2} & Distinguish between Intention and Competence with respect Trust with
  suitable example & 5 & CO3 & L2 & \ap{Q32} \\ \hline
\qn{3} & The complete content of respect is to be able to see that `the other
  is similar to me and we are complementary'. Define complementary in this
  context. & 5 & CO1 & L1 & \ap{Q67} \\ \hline
\qn{4} & Disrespect arises out of differentiation leading to discrimination on
  the basis of body, physical facility or beliefs. illustrate With an example
  & 5 & CO2 & L2 & \ap{Q36} \\ \hline
\qn{5} & Illustrate the natural acceptance values in relationship with an
  example for each (i)~Reverence (ii)~Glory & 5 & CO3 & L1 & \ap{Q39}
  \\ \hline
\end{qtable}

\subsection*{Part B --- marks distribution as printed}

\begin{mdtable}
Test --- Max Marks & 10 & 10 & 5 & --- & 15 & 10 & --- & --- & --- & --- \\
\hline
\end{mdtable}

\section{Improvement CIE --- June 2025}

\begin{papermeta}
\mk{Course code}   & \mv{HS248AT}                    &
\mk{Maximum marks} & \mv{30}                          \\
\mk{Semester}      & \mv{IV}                          &
\mk{Duration}      & \mv{60 minutes}                  \\
\mk{Academic year} & \mv{2024--2025, Even Semester}   &
\mk{Assessment}    & \mv{Improvement CIE --- Quiz and Test} \\
\mk{Date}          & \mv{June 2025}                   &
\mk{Structure}     & \mv{Quiz 5 M \textbf{+} Test 25 M}  \\
\mk{Set by}        & \mv{Department of Mechanical Engineering} &
                   & \\
\end{papermeta}

\subsection*{Quiz \hfill 5 Marks}

\begin{qtable}
\qn{1} & The participation of the human being in ensuring the role of physical
  facility in nurture, protection and providing means for the body is called
  its \blank
  \optsinline{Utility value}{Artistic value}{Activity}{Energy}
  & 1 & CO1 & L2 & \ap{A22} \\ \hline
\qn{2} & When something is active or has activity, it is called as \blank
  \optsinline{space}{unit}{energy}{value}
  & 1 & CO2 & L3 & \ap{A23} \\ \hline
\qn{3} & The following are the characteristics of space except \blank
  \optsinline{Unlimited}{No activity}{Energy in equilibrium}{Self-organized}
  & 1 & CO1 & L1 & \ap{A24} \\ \hline
\qn{4} & The natural characteristic of material order is \blank
  \optsinline{Composition}{Decomposition}{Composition/Decomposition}%
  {Nurture/worsen}
  & 1 & CO1 & L1 & \ap{A25} \\ \hline
\qn{5} & Conformance of material order is named as \blank
  \optsinline{constitution conformance}{Seed conformance}{Breed conformance}%
  {sanskaar conformance}
  & 1 & CO3 & L3 & \ap{A26} \\ \hline
\end{qtable}

\subsection*{Test \hfill 25 Marks}

\begin{qtable}
\qn{1} & Explain the classification of units into Four Orders and the harmony
  among them with the help of a schematic diagram. & 6 & CO1 & L2 &
  \ap{Q74} \\ \hline
\qn{2} & What do you mean by `conformance'? Explain the conformance in the
  four orders. & 6 & CO2 & L4 & \ap{Q55} \\ \hline
\qn{3} & Explain how there is recyclability and self-regulation in nature.
  & 6 & CO3 & L3 & \ap{Q60} \\ \hline
\qn{4} & Elaborate the terms `units' and `space' and mention their
  characteristics. & 7 & CO4 & L4 & \ap{Q65} \\ \hline
\end{qtable}

\subsection*{Marks distribution as printed}

\begin{mdtable}
Test --- Marks & 9 & 7 & 7 & 7 & 3 & 7 & 7 & 13 & --- & --- \\ \hline
\end{mdtable}

\vspace{4pt}
Both rows total 30, which is the paper's maximum. Note that L4 alone carries
13 of the 30 marks: this is an analytical paper, and the six- and seven-mark
questions on conformance and on units and space are where those marks sit.

\section{Semester End Examination --- June/July 2025}
\sectionmark{SEE --- June/July 2025}

\begin{papermeta}
\mk{Course code}   & \mv{HS248AT}                    &
\mk{Maximum marks} & \mv{50}                          \\
\mk{Semester}      & \mv{IV}                          &
\mk{Duration}      & \mv{2 hours}                     \\
\mk{Examination}   & \mv{IV Semester B.E. Regular / Supplementary} &
\mk{Assessment}    & \mv{Semester End Examination}    \\
\mk{Date}          & \mv{June/July 2025}              &
\mk{Applies to}    & \mv{Common to all programs}      \\
\end{papermeta}

\begin{note}
\textbf{\sffamily Instructions to the students, as printed.}
\begin{enumerate}[leftmargin=1.6em,itemsep=2pt,topsep=4pt]
\item Answer all questions from Part A. Part A questions should be answered in
      first three pages of the answer book only.
\item Answer THREE full questions from Part B. In Part B question number 2 is
      compulsory. Answer any one full question from 3 and 4 \& 5 and 6.
\end{enumerate}
\end{note}

The blueprint is worth reading off the instructions before the questions.
Part~A is ten compulsory one-mark objective questions. Part~B sets
\textbf{Q2 compulsory} (14 marks), then \textbf{one of Q3 or Q4} (13 marks) and
\textbf{one of Q5 or Q6} (13 marks) --- three full questions, 40 marks, and two
genuine choices to make in the hall.

\subsection*{Part A --- Compulsory Objective Questions \hfill 10 Marks}

\begin{qtable}
\qn{1.1} & How does ensuring justice in relationships contribute to society?
  \begin{optlist}
  \item It leads to fearlessness and co-existence
  \item It leads to prosperity and right understanding
  \item It promotes self-regulation and health
  \item It ensures better exchange and storage
  \end{optlist}
  & 01 & CO3 & L1 & \ap{A27} \\ \hline
\qn{1.2} & Imagination is continuous with
  \optsinline{Time}{Soul}{Body}{None}
  & 01 & CO1 & L1 & \ap{A28} \\ \hline
\qn{1.3} & Identify the solution which helps human being to transform from
  animal consciousness to human consciousness.
  \optsinline{Right understanding}{Realization}{Value education}%
  {Physical facilities.}
  & 01 & CO1 & L1 & \ap{A29} \\ \hline
\qn{1.4} & Which of the following is a fundamental human aspiration related to
  relationships?
  \begin{optlist}
  \item Physical comfort.
  \item Financial security.
  \item Mutual understanding and fulfillment.
  \item Accumulating wealth.
  \end{optlist}
  & 01 & CO2 & L2 & \ap{A30} \\ \hline
\qn{1.5} & Which of the following capacity leads to desires
  \optsinline{Power}{Expectation}{Realization}{Thoughts}
  & 01 & CO2 & L1 & \ap{A31} \\ \hline
\qn{1.6} & Which of the following is considered the foundation of a strong
  human relationship?
  \optsinline{Intelligence}{Trust}{Wealth}{Appearance}
  & 01 & CO1 & L2 & \ap{A32} \\ \hline
\qn{1.7} & Respect in a relationship means:
  \begin{optlist}
  \item Agreeing to everything
  \item Controlling the other person's actions
  \item Valuing the other person's opinions and boundaries
  \item Always trying to win argument.
  \end{optlist}
  & 01 & CO3 & L1 & \ap{A33} \\ \hline
\qn{1.8} & The participation of the human being in ensuring the role of
  physical facility to help and preserve its utility is called its \blank
  \optsinline{Utility value}{Artistic value}{Activity}{Energy}
  & 01 & CO3 & L2 & \ap{A34} \\ \hline
\qn{1.9} & A harmonious world is created by values at 4 levels. These are:
  \begin{optlist}
  \item Home, family, society, country
  \item Individual, family, society, universe
  \item School, home, office, temple
  \item None of these
  \end{optlist}
  & 01 & CO1 & L1 & \ap{A35} \\ \hline
\qn{1.10} & When nature is submerged in space it is known as
  \optsinline{Conformance}{Acceptance}{Mixing}{Co-existence}
  & 01 & CO3 & L2 & \ap{A36} \\ \hline
\end{qtable}

\subsection*{Part B --- Descriptive Questions \hfill 40 Marks}

\vspace{2pt}
\textbf{\sffamily\color{slate}Question 2 --- compulsory}

\begin{qtable}
\qn{2a} & What do you understand by the value of an entity? What is the value
  of a human being? & 08 & CO1 & L2 & \ap{Q4} \\ \hline
\qn{2b} & For human being physical facility is necessary, but relationship is
  also necessary: Justify with an example comparing with animal order.
  & 06 & CO2 & L4 & \ap{Q21} \\ \hline
\end{qtable}

\vspace{4pt}
\textbf{\sffamily\color{slate}Questions 3 and 4 --- answer any one full
question}

\begin{qtable}
\qn{3a} & ``Discuss how right understanding in relationships contributes to
  harmony in the family. Analyze the role of trust, respect, and
  responsibility in maintaining family unity, and provide examples to support
  your answer.'' & 07 & CO4 & L4 & \ap{Q68} \\ \hline
\qn{3b} & Articulate the four goals of human beings living in a society.
  Critically examine the state of the society today in context with the
  fulfilment of a comprehensive human goal. & 06 & CO3 & L3 &
  \ap{Q41}\newline\ap{Q43} \\ \hline
\orband
\qn{4a} & How do justice and preservation contribute to a harmonious society?
  Analyze the shortcomings of focusing solely on punishment in cases of
  injustice, and explain how true preservation ensures prosperity and
  coexistence with nature. & 06 & CO3 & L4 & \ap{Q69} \\ \hline
\qn{4b} & Explain the concept of harmony in the family. Why is it important,
  and how can it be achieved through right understanding and mutual
  relationships? & 07 & CO3 & L2 & \ap{Q70} \\ \hline
\end{qtable}

\vspace{4pt}
\textbf{\sffamily\color{slate}Questions 5 and 6 --- answer any one full
question}

\begin{qtable}
\qn{5a} & What is sanskaar? Explain its effects or the conformance of the
  human order. & 06 & CO4 & L4 & \ap{Q55} \\ \hline
\qn{5b} & What is the svabhava (natural characteristic) of a unit? Mention the
  natural characteristics of the material, pranic, animal and human orders.
  Give examples & 07 & CO4 & L3 & \ap{Q54} \\ \hline
\orband
\qn{6a} & What do you mean by `conformance'? Explain the conformance in the
  four orders. & 07 & CO3 & L4 & \ap{Q55} \\ \hline
\qn{6b} & Explain the meaning of co-existence. & 06 & CO3 & L3 & \ap{Q63}
  \\ \hline
\end{qtable}

\section{Improvement CIE --- 18 June 2026}

\begin{papermeta}
\mk{Course code}   & \mv{HS248AT}                    &
\mk{Maximum marks} & \mv{05 $+$ 25}                   \\
\mk{Semester}      & \mv{IV}                          &
\mk{Duration}      & \mv{10 min $+$ 60 min}           \\
\mk{Date}          & \mv{18 June 2026}                &
\mk{Assessment}    & \mv{Improvement CIE --- Quiz $+$ Test} \\
\mk{Structure}     & \mv{Part A 5 M \textbf{+} Part B 25 M} &
\mk{Format}        & \mv{Short answer, not multiple choice} \\
\end{papermeta}

This paper breaks the pattern of the earlier improvement CIEs in one respect
worth noting: \textbf{Part~A is not multiple choice}. It sets five one-mark
short-answer questions to be written in ten minutes --- definitions, lists and
one-line differentiations. There is nothing to guess at.

\subsection*{Part A --- Short Answer Questions \hfill 5 Marks}

\begin{qtable}
\qn{1} & Give any two definitions for the term co-existence.
  & 1 & CO1 & L1 & \ap{C1} \\ \hline
\qn{2} & State the four orders of nature. & 1 & CO2 & L1 & \ap{C2} \\ \hline
\qn{3} & Identify the svabhava\,/\,value of the self (`I') in human beings.
  & 1 & CO2 & L2 & \ap{C3} \\ \hline
\qn{4} & Summarize 2 things that human beings have to do to realize the mutual
  fulfilment in nature. & 1 & CO3 & L2 & \ap{C4} \\ \hline
\qn{5} & Differentiate between Units and Space. & 1 & CO4 & L2 & \ap{C5}
  \\ \hline
\end{qtable}

\subsection*{Part B --- Descriptive Questions \hfill 25 Marks}

\begin{qtable}
\qn{1} & Describe how harmony is achieved among the four orders of nature.
  & 5 & CO4 & L2 & \ap{Q59}\newline\ap{Q74} \\ \hline
\qn{2} & Justify that requirement of any order is already available in
  abundance in nature. & 5 & CO3 & L4 & \ap{Q72} \\ \hline
\qn{3} & Demonstrate how interconnectedness, self-regulation and mutual
  fulfilment are among the four orders of nature with an example.
  & 5 & CO2 & L3 & \ap{Q59}\newline\ap{Q60} \\ \hline
\qn{4} & Summarize how the domain of consciousness helps in accomplishing
  development or permanent change. & 5 & CO4 & L2 & \ap{Q73} \\ \hline
\qn{5} & Outline 2 types of units with examples. & 5 & CO3 & L2 & \ap{Q71}
  \\ \hline
\end{qtable}

\subsection*{Marks distribution as printed}

\begin{mdtable}
Test --- Max Marks & 1 & 7 & 11 & 11 & 2 & 18 & 5 & 5 & --- & --- \\ \hline
\end{mdtable}

\vspace{4pt}
Both rows total 30. L2 carries 18 of the 30 marks, which is the mirror image of
the June 2025 paper: this one rewards clear explanation rather than analysis.
